\documentclass[10pt,aspectratio=169]{beamer}

\usepackage[utf8]{inputenc}
\usepackage[T1]{fontenc}
\usepackage[french]{babel}
\usepackage{amsmath,amssymb,booktabs,array,multirow}
\usepackage{graphicx}
\usepackage{tikz}
\usepackage{comment}
\usepackage{xcolor}
\usepackage{colortbl}
\usetikzlibrary{arrows.meta,positioning,calc,shadows,decorations.pathreplacing,shapes.geometric}
\graphicspath{{assets/}}

% ---------- Theme ----------
\usetheme{Madrid}
\usecolortheme{default}
\usefonttheme{serif}
\setbeamertemplate{navigation symbols}{}
\setbeamertemplate{blocks}[rounded][shadow=true]

% ---------- Colors ----------
\definecolor{myblue}{RGB}{0,70,140}
\definecolor{myorange}{RGB}{220,100,0}
\definecolor{mygreen}{RGB}{0,130,60}
\definecolor{mygray}{RGB}{80,80,80}
\definecolor{rougebox}{RGB}{180,30,30}
\definecolor{bleubox}{RGB}{30,80,160}
\definecolor{vertbox}{RGB}{0,128,0}
\definecolor{orangeiter}{RGB}{210,120,0}
\definecolor{bleutitre}{RGB}{31,73,125}

\setbeamercolor{structure}{fg=myblue}
\setbeamercolor{frametitle}{bg=myblue, fg=white}
\setbeamercolor{title}{fg=white}
\setbeamercolor{block title}{bg=myblue!80, fg=white}
\setbeamercolor{block body}{bg=myblue!8, fg=black}
\setbeamercolor{alertblock title}{bg=red!70!black, fg=white}
\setbeamercolor{alertblock body}{bg=red!8, fg=black}
\setbeamercolor{exampleblock title}{bg=mygreen!80, fg=white}
\setbeamercolor{exampleblock body}{bg=mygreen!8, fg=black}

\setbeamertemplate{footline}{%
  \leavevmode\hbox{%
    \begin{beamercolorbox}[wd=.33\paperwidth,ht=2.4ex,dp=1ex,center]{author in head/foot}%
      \usebeamerfont{author in head/foot}\insertshortauthor
    \end{beamercolorbox}%
    \begin{beamercolorbox}[wd=.34\paperwidth,ht=2.4ex,dp=1ex,center]{title in head/foot}%
      \usebeamerfont{title in head/foot}\insertshorttitle
    \end{beamercolorbox}%
    \begin{beamercolorbox}[wd=.33\paperwidth,ht=2.4ex,dp=1ex,right]{date in head/foot}%
      \usebeamerfont{date in head/foot}\insertframenumber{} / \inserttotalframenumber\hspace*{2ex}
    \end{beamercolorbox}}%
}

\newcommand{\figcol}[1]{\includegraphics[width=\linewidth,height=0.88\textheight,keepaspectratio]{#1}}
\newcommand{\figbig}[1]{\includegraphics[width=\linewidth,height=0.82\textheight,keepaspectratio]{#1}}
\newcommand{\figtall}[1]{\includegraphics[height=0.84\textheight,keepaspectratio]{#1}}

\title[Full Adder, activations et CORDIC]{Optimisation des fonctions d'activation pour les réseaux de neurones\\\large Étude appliquée au \textit{Full Adder} et implémentation CORDIC sur FPGA}
\author[Ahchouch \& Benlahsar]{Fatima-Ezzahra Ahchouch \quad Lina Benlahsar}
\date{Mars 2026}

\begin{document}

% ============================================================
\begin{frame}[plain]
  \titlepage
\end{frame}

% ============================================================
% ============================================================
\begin{frame}{Sommaire}
  \tableofcontents
\end{frame}

% ============================================================
\section{Le Full Adder comme base de modélisation}

\begin{frame}{Le Full Adder : table, équations et décomposition logique}
\begin{columns}[T]
\begin{column}{0.39\textwidth}
  \begin{block}{Table de vérité}
    \centering\small
    \renewcommand{\arraystretch}{1.05}
    \begin{tabular}{ccccc}
      \toprule
      \(A\) & \(B\) & \(C_{in}\) & \(S\) & \(C_{out}\) \\
      \midrule
      0 & 0 & 0 & 0 & 0 \\
      0 & 0 & 1 & 1 & 0 \\
      0 & 1 & 0 & 1 & 0 \\
      0 & 1 & 1 & 0 & 1 \\
      1 & 0 & 0 & 1 & 0 \\
      1 & 0 & 1 & 0 & 1 \\
      1 & 1 & 0 & 0 & 1 \\
      1 & 1 & 1 & 1 & 1 \\
      \bottomrule
    \end{tabular}
  \end{block}
\end{column}
\begin{column}{0.57\textwidth}
  \begin{block}{Équations logiques}
    \[
      S = A \oplus B \oplus C_{in}
      \qquad
      C_{out}=AB + AC_{in}+BC_{in}
    \]
  \end{block}
  \centering
  \begin{tikzpicture}[>=Stealth, node distance=0.8cm,
    box/.style={draw,rounded corners,thick,minimum width=1.6cm,minimum height=0.55cm,align=center,fill=myblue!8},
    gate/.style={draw,rounded corners,thick,minimum width=1.45cm,minimum height=0.65cm,align=center,fill=mygreen!10}]
    \node[box] (A) at (0,1.2) {\(A\)};
    \node[box] (B) at (0,0.0) {\(B\)};
    \node[box] (C) at (0,-1.2) {\(C_{in}\)};
    \node[gate,fill=rougebox!10,draw=rougebox] (xor1) at (2.4,0.6) {XOR};
    \node[gate,fill=rougebox!10,draw=rougebox] (xor2) at (4.8,0.6) {XOR};
    \node[gate,fill=vertbox!10,draw=vertbox] (maj) at (3.7,-1.0) {Majorité};
    \node[box,fill=rougebox!10,draw=rougebox] (S) at (7.0,0.6) {\(S\)};
    \node[box,fill=vertbox!10,draw=vertbox] (CO) at (7.0,-1.0) {\(C_{out}\)};
    \draw[->,thick] (A)--(xor1); \draw[->,thick] (B)--(xor1);
    \draw[->,thick] (xor1)--(xor2); \draw[->,thick] (C)--(xor2);
    \draw[->,thick] (xor2)--(S);
    \draw[->,thick] (A.east)--++(0.75,0)--++(0,-1.65)--(maj.west);
    \draw[->,thick] (B.east)--++(0.45,0)--++(0,-1.0)--(maj.west);
    \draw[->,thick] (C.east)--++(0.55,0)--(maj.west);
    \draw[->,thick] (maj)--(CO);
  \end{tikzpicture}
 
  
\end{column}
\end{columns}
\begin{exampleblock}{Lecture mathématique}
    \small
    La sortie \(S\) contient la difficulté du XOR non séparable, tandis que \(C_{out}\) est une fonction seuil de majorité : les deux sorties exigent donc une architecture capable de combiner décision binaire et non-linéarité.
  \end{exampleblock}
\end{frame}

% ============================================================
\section{Neurone formel et apprentissage}

\begin{frame}{Structure d'un neurone artificiel appliquée au Full Adder}
\begin{columns}[T]
\begin{column}{0.61\textwidth}
\begin{center}
\begin{tikzpicture}[scale=0.92,
  input/.style={circle, fill=myblue!15, draw=myblue, thick, minimum size=0.68cm, font=\footnotesize},
  op/.style={circle, fill=myorange!20, draw=myorange, thick, minimum size=1.0cm, font=\large},
  act/.style={rectangle, rounded corners=4pt, fill=mygreen!18, draw=mygreen, thick, minimum width=1.18cm, minimum height=0.70cm, font=\footnotesize},
  arr/.style={-Stealth, thick, gray}
]
  \node[input] (a) at (0,1.45) {$A$};
  \node[input] (b) at (0,0.45) {$B$};
  \node[input] (c) at (0,-0.55) {$C_{in}$};
  \node[input] (bias) at (0,-1.55) {$1$};

  \node[op, drop shadow] (sigma) at (3.05,0) {$\Sigma$};
  \node[act, drop shadow] (f) at (5.05,0) {$f(z)$};
  \node[act, drop shadow,fill=rougebox!15,draw=rougebox] (out) at (7.0,0.58) {$S$};
  \node[act, drop shadow,fill=vertbox!15,draw=vertbox] (cout) at (7.0,-0.58) {$C_{out}$};

  \node[font=\footnotesize\bfseries,myorange] at (3.05,1.12) {somme pondérée};
  \node[font=\footnotesize\bfseries,mygreen!60!black] at (5.05,0.75) {activation};
  \node[font=\footnotesize\bfseries,rougebox] at (8.1,0.58) {parité};
  \node[font=\footnotesize\bfseries,vertbox] at (8.25,-0.58) {majorité};

  \draw[arr] (a) -- node[above,sloped,font=\tiny]{$w_A$} (sigma);
  \draw[arr] (b) -- node[above,sloped,font=\tiny]{$w_B$} (sigma);
  \draw[arr] (c) -- node[below,sloped,font=\tiny]{$w_C$} (sigma);
  \draw[arr,dashed] (bias) -- node[below,sloped,font=\tiny]{$b$} (sigma);
  \draw[arr] (sigma) -- node[above,font=\tiny]{$z$} (f);
  \draw[arr] (f) -- (out);
  \draw[arr] (f) -- (cout);
\end{tikzpicture}
\end{center}

\vspace{-0.1cm}
\begin{block}{Modèle mathématique}
\[
  z = w_A A + w_B B + w_C C_{in} + b,\qquad y=f(z)
\]
\end{block}
\end{column}

\begin{column}{0.36\textwidth}
\begin{exampleblock}{Commentaire}
\footnotesize
Un neurone isolé réalise une séparation par seuil. Il modélise naturellement une décision simple ou une majorité partielle, mais pas la parité du bit somme.
\end{exampleblock}
\vspace{0.10cm}
\begin{alertblock}{Conséquence}
\footnotesize
La sortie \(S=A\oplus B\oplus C_{in}\) nécessite des couches cachées et une activation non linéaire ; \(C_{out}\) est plus proche d'une fonction de seuil.
\end{alertblock}
\vspace{0.05cm}
\begin{center}
\begin{tikzpicture}[scale=0.76]
\draw[->] (-1.7,0)--(1.8,0) node[right, font=\tiny]{\(z\)};
\draw[->] (0,-1.1)--(0,1.15) node[above, font=\tiny]{\(f(z)\)};
\draw[mygreen,thick,domain=-1.6:1.6,samples=50] plot (\x,{(exp(\x)-exp(-\x))/(exp(\x)+exp(-\x))});
\draw[rougebox,thick] (-1.6,-0.85)--(0,-0.85);
\draw[rougebox,thick] (0,0.85)--(1.6,0.85);
\node[mygreen,font=\tiny] at (1.1,0.55){\(\tanh\)};
\node[rougebox,font=\tiny] at (1.1,0.95){step};
\end{tikzpicture}
\end{center}
\end{column}
\end{columns}
\end{frame}

\begin{frame}{Propagation avant : calcul de \((\hat S,\hat C_{out})\)}
\centering
\begin{tikzpicture}[>=Stealth,
  layer/.style={draw,rounded corners,fill=myblue!8,minimum width=1.55cm,minimum height=0.65cm,align=center,font=\scriptsize},
  act/.style={draw,rounded corners,fill=mygreen!10,minimum width=1.55cm,minimum height=0.65cm,align=center,font=\scriptsize}]
  \node[layer] (in) at (0,0) {Entrées\\\(A,B,C_{in}\)};
  \node[act] (h1) at (2.2,0.8) {Couche 1\\\(\tanh\)};
  \node[act] (h2) at (4.4,0.8) {Couche 2\\\(\tanh\)};
  \node[act] (h3) at (6.6,0.8) {Couche 3\\\(\tanh\)};
  \node[layer,fill=rougebox!10,draw=rougebox] (s) at (8.8,1.25) {\(\hat S\)};
  \node[layer,fill=vertbox!10,draw=vertbox] (co) at (8.8,0.25) {\(\hat C_{out}\)};
  \draw[->,thick,myblue] (in)--(h1); \draw[->,thick,myblue] (h1)--(h2); \draw[->,thick,myblue] (h2)--(h3); \draw[->,thick,myblue] (h3)--(s); \draw[->,thick,myblue] (h3)--(co);
  \node[draw,myorange,fill=myorange!8,rounded corners,align=center,font=\scriptsize] at (4.4,-1.0) {Accumulation pondérée : \(z=W x+b\)};
\end{tikzpicture}
\begin{columns}[T]
\begin{column}{0.45\textwidth}
  \begin{block}{Calcul direct}
    \small
    La propagation avant applique successivement les poids, les biais et les activations afin d'estimer les deux sorties :
    \[
      (A,B,C_{in}) \longrightarrow h^{(1)} \longrightarrow h^{(2)} \longrightarrow (\hat S,\hat C_{out}).
    \]
  \end{block}
 
\end{column}
\begin{column}{0.52\textwidth}
 \begin{exampleblock}{Objectif de la propagation}
    \small
    Chaque ligne de la table du Full Adder doit être projetée vers les deux bits logiques attendus.
  \end{exampleblock}
\end{column}
\end{columns}
\end{frame}

\begin{frame}{Rétropropagation : correction des poids}
\begin{columns}[T]
\begin{column}{0.48\textwidth}
  \begin{block}{Erreur globale}
    \[
      E = \frac{1}{2}\big[(S-\hat S)^2 + (C_{out}-\hat C_{out})^2\big]
    \]
  \end{block}
  \begin{block}{Mise à jour}
    \[
      w \leftarrow w - \eta\,\frac{\partial E}{\partial w},
      \qquad
      b \leftarrow b - \eta\,\frac{\partial E}{\partial b}
    \]
  \end{block}
  \begin{alertblock}{Condition essentielle}
    \small
    La rétropropagation nécessite une activation différentiable dans les couches cachées : \(\tanh\) ou sigmoïde. La fonction step est utilisée en sortie ou pour une construction logique exacte.
  \end{alertblock}
\end{column}
\begin{column}{0.48\textwidth}
\centering
\begin{tikzpicture}[
  box/.style={rectangle, rounded corners, draw=myblue, fill=myblue!8, thick, minimum width=2.6cm, minimum height=0.75cm, align=center, font=\scriptsize},
  ebox/.style={rectangle, rounded corners, draw=red!70!black, fill=red!8, thick, minimum width=2.6cm, minimum height=0.75cm, align=center, font=\scriptsize},
  gbox/.style={rectangle, rounded corners, draw=myorange, fill=myorange!8, thick, minimum width=2.6cm, minimum height=0.75cm, align=center, font=\scriptsize},
  fwd/.style={-Stealth, thick, myblue}, bwd/.style={-Stealth, thick, myorange}
]
  \node[box]  (inp)  at (0,0) {Entrées\\\(A,B,C_{in}\)};
  \node[box]  (pred) at (4.1,0) {Sorties\\\(\hat S,\hat C_{out}\)};
  \node[ebox] (err)  at (4.1,-1.7) {Erreur\\\(E\)};
  \node[gbox] (grad) at (0,-1.7) {Gradients\\\(\nabla E\)};
  \draw[fwd] (inp)--node[above,font=\tiny]{propagation}(pred);
  \draw[fwd] (pred)--node[right,font=\tiny]{comparaison}(err);
  \draw[bwd] (err)--node[below,font=\tiny]{rétropropagation}(grad);
  \draw[bwd] (grad)--node[left,font=\tiny]{mise à jour}(inp);
\end{tikzpicture}
\vspace{0.3cm}
\begin{exampleblock}{Lecture Full Adder}
\small
Les couches cachées ajustent leurs frontières internes afin de reproduire simultanément la parité \(S\) et la majorité \(C_{out}\).
\end{exampleblock}
\end{column}
\end{columns}
\end{frame}

% ============================================================
\section{Fonctions d'activation et cas logiques}

\begin{frame}{Pourquoi une fonction d'activation ?}
\begin{center}
\begin{tikzpicture}[scale=1.30]
  \node[font=\bfseries\footnotesize,red!70!black] at (0.55,1.65) {AND : séparable linéairement};
  \begin{scope}[xshift=-1.1cm]
    \draw[->] (-0.2,0)--(1.45,0) node[right,font=\tiny]{\(x_1\)};
    \draw[->] (0,-0.2)--(0,1.45) node[above,font=\tiny]{\(x_2\)};
    \fill[red] (0,0) circle (2.4pt);
    \fill[red] (0,1) circle (2.4pt);
    \fill[red] (1,0) circle (2.4pt);
    \fill[mygreen] (1,1) circle (2.8pt);
    \draw[myblue,ultra thick] (-0.1,1.15)--(1.15,-0.1);
  \end{scope}
  \node[font=\bfseries\footnotesize,mygreen!70!black] at (5.15,1.65) {XOR : non séparable};
  \begin{scope}[xshift=3.6cm]
    \draw[->] (-0.2,0)--(1.45,0) node[right,font=\tiny]{\(x_1\)};
    \draw[->] (0,-0.2)--(0,1.45) node[above,font=\tiny]{\(x_2\)};
    \fill[red] (0,0) circle (2.4pt);
    \fill[mygreen] (0,1) circle (2.8pt);
    \fill[mygreen] (1,0) circle (2.8pt);
    \fill[red] (1,1) circle (2.4pt);
    \draw[mygreen,ultra thick,rounded corners=8pt] (-0.12,0.28)--(0.28,-0.12)--(1.12,0.72)--(0.72,1.12)--cycle;
    \draw[rougebox,dashed,thick] (0.5,-0.15)--(0.5,1.25);
  \end{scope}
\end{tikzpicture}
\end{center}
\vspace{-0.1cm}
\begin{columns}[T]
\begin{column}{0.32\textwidth}
  \begin{block}{Sans non-linéarité}
  \scriptsize
  Des couches linéaires empilées restent équivalentes à une seule transformation linéaire ; elles ne résolvent pas XOR.
  \end{block}
\end{column}
\begin{column}{0.32\textwidth}
  \begin{exampleblock}{Avec activation}
  \scriptsize
  Une activation comme \(\tanh\) courbe l'espace de représentation et permet des frontières de décision composées.
  \end{exampleblock}
\end{column}
\begin{column}{0.32\textwidth}
  \begin{alertblock}{Cas Full Adder}
  \scriptsize
  \(C_{out}\) relève d'une majorité, mais \(S=A\oplus B\oplus C_{in}\) impose non-linéarité et couches cachées.
  \end{alertblock}
\end{column}
\end{columns}
\end{frame}

\begin{frame}{Les fonctions d'activation cibles : step, sigmoïde et tanh}
\begin{columns}[T]
\begin{column}{0.30\textwidth}
  \begin{block}{Step}
    \[
      f(x)=\begin{cases}1 & x\ge 0\\0 & x<0\end{cases}
    \]
    \small
    \begin{itemize}
      \item Décision logique nette.
      \item Non différentiable.
      \item Adaptée au seuillage final.
    \end{itemize}
  \end{block}
\end{column}
\begin{column}{0.32\textwidth}
  \begin{block}{Sigmoïde}
    \[
      \sigma(x)=\frac{1}{1+e^{-x}},\quad 
    \]
    \small
    \begin{itemize}
      \item Sortie dans \((0,1)\).
      \item Gradient maximum \(0.25\).
      \item Coûteuse en exponentielle.
    \end{itemize}
  \end{block}
\end{column}
\begin{column}{0.28\textwidth}
  \begin{block}{Tangente hyperbolique}
    \[
      \tanh(x)=\frac{e^x-e^{-x}}{e^x+e^{-x}},\quad 
    \]
    \small
    \begin{itemize}
      \item Sortie dans \((-1,1)\).
      \item Centrée en zéro.
      \item Gradient maximum \(1\).
    \end{itemize}
  \end{block}
\end{column}
\end{columns}
\vspace{0.18cm}
\begin{exampleblock}{Relation fondamentale}
\[
  \tanh(x)=2\sigma(2x)-1
  \qquad \Longleftrightarrow \qquad
  \sigma(x)=\frac{1+\tanh(x/2)}{2}
\]
\end{exampleblock}
\end{frame}

% ============================================================
\section{Études préalables : ET et XOR}

\begin{frame}{Étude préalable : ET logique comme baseline de convergence}
\begin{columns}[T]
\begin{column}{0.28\textwidth}
  \begin{block}{Fonction linéairement séparable}
    \small
    \begin{itemize}
      \item Un neurone à seuil peut réaliser la fonction ET.
      \item Les frontières de décision sont simples.
      \item La convergence sert à valider la chaîne d'apprentissage.
    \end{itemize}
  \end{block}
  \begin{exampleblock}{Lien avec \(C_{out}\)}
    \small
    La retenue du Full Adder repose sur des produits logiques \(AB\), \(AC_{in}\), \(BC_{in}\) puis une majorité.
  \end{exampleblock}
\end{column}
\begin{column}{0.70\textwidth}

  \includegraphics[width=\linewidth,height=5.5cm]{et.png}
\end{column}
\end{columns}
\end{frame}

\begin{frame}{Étude de cas : XOR, cœur non linéaire de la somme}
\begin{columns}[T]
\begin{column}{0.28\textwidth}
  \begin{block}{Non-séparabilité}
    \small
    \begin{itemize}
      \item XOR ne peut pas être séparé par une seule droite.
      \item Une couche cachée crée des frontières intermédiaires.
      \item Le bit \(S\) du Full Adder est une parité : \(A\oplus B\oplus C_{in}\).
    \end{itemize}
  \end{block}
 
\end{column}
\begin{column}{0.70\textwidth}
  \figcol{xor.png}
\end{column}
\end{columns}
 \begin{exampleblock}{Conclusion expérimentale}
    \small
    Tanh converge plus rapidement que sigmoïde grâce à son gradient plus grand et à sa sortie centrée en zéro.
  \end{exampleblock}
\end{frame}

% ============================================================
\section{Réseau de neurones pour le Full Adder}

\begin{frame}{Architecture Full Adder : combinaison tanh + step}
\begin{center}
\includegraphics[width=0.94\linewidth,height=0.46\textheight,keepaspectratio]{adder_overview.png}
  \begin{block}{Architecture profonde proposée}
    \centering
    \fbox{\(3 \rightarrow 16 \rightarrow 12 \rightarrow 8 \rightarrow 2 \rightarrow \text{step}\)}\\[-1pt]
    {\footnotesize entrées \(A,B,C_{in}\) \(\rightarrow\) couches cachées \(\tanh\) \(\rightarrow\) sorties continues \(\rightarrow\) décision binaire}
  \end{block}
  
\end{center}
\vspace{-0.08cm}
\end{frame}
\begin{frame}{Architecture Full Adder : combinaison tanh + step}
\begin{center}
\includegraphics[width=0.94\linewidth,height=0.46\textheight,keepaspectratio]{adder_overview.png}
\end{center}
\begin{columns}[T]
\begin{column}{0.49\textwidth}
  \begin{exampleblock}{Pourquoi plusieurs couches cachées ?}
  \scriptsize
  La parité à trois entrées est plus complexe qu'un XOR simple. La profondeur factorise les combinaisons élémentaires, les motifs XOR puis la fusion vers \(S\) et \(C_{out}\).
  \end{exampleblock}
\end{column}
\begin{column}{0.49\textwidth}
  \begin{alertblock}{Rôle des activations}
  \scriptsize
  \textbf{Tanh} assure l'apprentissage différentiable dans les couches cachées. \textbf{Step} transforme les deux sorties continues en bits logiques \(0/1\).
  \end{alertblock}
\end{column}
\end{columns}
\end{frame}
\begin{comment}
    

\begin{frame}{Full Adder : réseau step exact comme référence logique}
\begin{columns}[T]
\begin{column}{0.45\textwidth}
  \begin{block}{Architecture constructive}
    \centering
    \fbox{\(3\rightarrow 8\rightarrow 2\)}
  \end{block}
  \begin{itemize}\small
    \item Les 8 neurones cachés détectent les 8 mintermes de la table de vérité.
    \item Les sorties réalisent un OR des mintermes associés à \(S=1\) et \(C_{out}=1\).
    \item Aucun apprentissage n'est nécessaire : les poids peuvent être fixés analytiquement.
  \end{itemize}
  \begin{exampleblock}{Intérêt}
    \small
    Cette architecture donne une référence exacte pour valider les résultats du réseau entraînable.
  \end{exampleblock}
\end{column}
\begin{column}{0.55\textwidth}

  \includegraphics[width=\linewidth,height=5.5cm]{adder_step.png}
\end{column}
\end{columns}
\end{frame}
\end{comment}
\begin{frame}{Full Adder : réseau tanh profond puis seuillage step}
\begin{columns}[T]
\begin{column}{0.45\textwidth}
  \begin{block}{Architecture entraînable}
    \centering
    \fbox{\(3\rightarrow 16\rightarrow 12\rightarrow 8\rightarrow 2\)}
  \end{block}
  \begin{itemize}\small
    \item Tanh dans toutes les couches cachées.
    \item Sorties continues converties par step : \(y\ge 0.5\Rightarrow 1\), sinon \(0\).
    \item Optimisation par rétropropagation sur les 8 lignes du Full Adder.
  \end{itemize}
  \begin{alertblock}{Lien vers le FPGA}
    \small
    La présence de \(\tanh\) rend nécessaire une implémentation matérielle efficace de cette activation.
  \end{alertblock}
\end{column}
\begin{column}{0.55\textwidth}
  
   \includegraphics[width=\linewidth,height=5cm]{adder_tanh.png}
\end{column}
\end{columns}
\end{frame}
\begin{comment}
    

\begin{frame}{Full Adder : synthèse des deux approches}
\begin{columns}[T]
\begin{column}{0.50\textwidth}
  \includegraphics[width=\linewidth,height=5.5cm]{adder_overview.png}
\end{column}
\begin{column}{0.50\textwidth}
  \begin{block}{Comparaison fonctionnelle}
    \small
    \renewcommand{\arraystretch}{1.2}
    \begin{tabular}{p{2cm}p{2cm}p{2cm}}
      \toprule
      & \textbf{Step exact} & \textbf{Tanh + step} \\
      \midrule
      Nature & constructive & entraînable \\
      Couches cachées & 1 & 3 \\
      Activation cachée & step & tanh \\
      Sortie & step & step \\
      Usage & référence logique & réseau différentiable \\
      \bottomrule
    \end{tabular}
  \end{block}
 \begin{exampleblock}{Transition}
    \small
    L'architecture tanh + step motive directement l'étude de l'approximation de \(\tanh\) en matériel numérique.
  \end{exampleblock}
\end{column}
\end{columns}
 
\end{frame}
\end{comment}

% ============================================================
\section{Implémentation matérielle des activations}

\begin{frame}{Problématique matérielle : pourquoi optimiser \(\tanh\) et \(\sigma\) ?}
\begin{columns}[T]
\begin{column}{0.48\textwidth}
  \begin{alertblock}{Calcul direct coûteux sur FPGA}
    \small
    \begin{itemize}
      \item \(\sigma(x)=\dfrac{1}{1+e^{-x}}\) nécessite exponentielle et division.
      \item \(\tanh(x)=\dfrac{e^x-e^{-x}}{e^x+e^{-x}}\) nécessite aussi des opérations coûteuses.
      \item Les blocs DSP et les ressources logiques deviennent rapidement limitants.
    \end{itemize}
  \end{alertblock}
\end{column}
\begin{column}{0.48\textwidth}
  \begin{block}{Trois approches étudiées}
    \begin{enumerate}\small
      \item \textbf{LUT} : table de correspondance.
      \item \textbf{Taylor} : approximation polynomiale.
      \item \textbf{CORDIC} : additions et décalages binaires.
    \end{enumerate}
  \end{block}
  \begin{exampleblock}{Critères}
    \small
    Précision, stabilité sur l'intervalle utile, faible mémoire, absence de multiplieurs DSP.
  \end{exampleblock}
\end{column}
\end{columns}
\vspace{0.25cm}
\centering
\begin{tikzpicture}
\draw[ultra thick,myblue,-Stealth] (0,0) -- (10,0)
node[midway,above,font=\small\bfseries]{Objectif : remplacer les exponentielles par une architecture shifts + additions};
\end{tikzpicture}
\end{frame}

% ============================================================
\section{Méthodes classiques d'approximation}

\begin{frame}{Méthode LUT : simple, mais pas forcément optimale}
\begin{columns}[T]
\begin{column}{0.28\textwidth}
  \begin{block}{Principe}
    \small
    La fonction est échantillonnée puis stockée en ROM :
    \[
      x_i=x_{min}+i\Delta x,
      \\
      \qquad f(x)\approx \mathrm{LUT}[i].
    \]
  \end{block}
  \begin{itemize}\small
    \item Sans interpolation : erreur en dents de scie.
    \item Avec interpolation : précision améliorée.
    \item Interpolation : coût d'un multiplieur.
  \end{itemize}
\end{column}
\begin{column}{0.70\textwidth}
  \figcol{lut.png}
\end{column}
\end{columns}
\end{frame}

\begin{frame}{Développement de Taylor : précis localement, fragile globalement}
\begin{columns}[T]
\begin{column}{0.28\textwidth}
  \begin{block}{Approximation polynomiale}
    \small
    \[
      P_3(x)=x-\frac{x^3}{3}
    \]
    \[
      P_5(x)=x-\frac{x^3}{3}+\frac{2x^5}{15}
    \]
  \end{block}
  \begin{alertblock}{Limites}
    \small
    \begin{itemize}
      \item Précision locale autour de zéro.
      \item Divergence pour les grands \(|x|\).
      \item Multiplications coûteuses en FPGA.
    \end{itemize}
  \end{alertblock}
\end{column}
\begin{column}{0.70\textwidth}

  \includegraphics[width=\linewidth,height=5.5cm]{taylor.png}
\end{column}
\end{columns}
\end{frame}

\begin{frame}{Comparaison des méthodes d'approximation sur \(\tanh\)}
\begin{columns}[T]
\begin{column}{0.28\textwidth}
  \begin{block}{Lecture des erreurs}
    \small
    \begin{itemize}
      \item LUT : stable mais quantifiée.
      \item Taylor : bonne près de 0, instable hors zone locale.
      \item CORDIC : erreur très faible et plus uniforme.
    \end{itemize}
  \end{block}
  \begin{alertblock}{Conclusion}
    \small
    Une solution matérielle robuste doit éviter à la fois la mémoire excessive, les multiplications et la divergence hors domaine.
  \end{alertblock}
\end{column}
\begin{column}{0.70\textwidth}

  \includegraphics[width=\linewidth,height=5.5cm]{bilan.png}
\end{column}
\end{columns}
\end{frame}

\begin{frame}{Bilan comparatif : choix de la méthode CORDIC}
\centering
\renewcommand{\arraystretch}{1.55}
\begin{tabular}{|l|c|c|c|p{4.0cm}|}
\hline
\rowcolor{myblue!20}
\textbf{Méthode} & \textbf{MSE} & \textbf{Stabilité} & \textbf{DSP ?} & \textbf{Analyse} \\
\hline
LUT \(N=32\) sans interp. & \(\approx 10^{-3}\) & Stable & Non & Simple, mais effet escalier \\
\hline
LUT \(N=16\) avec interp. & \(\approx 2\times10^{-5}\) & Stable & Oui & Précise, mais nécessite un multiplieur \\
\hline
Taylor \(n=7\) & \(\approx 10^{-4}\) & Diverge & Oui & Peu robuste hors voisinage de zéro \\
\hline
\rowcolor{mygreen!20}
\textbf{CORDIC \(k=16\)} & \(\mathbf{\approx 10^{-8}}\) & \textbf{Uniforme} & \textbf{Non} & \textbf{Solution retenue pour \(\tanh\) et \(\sigma\)} \\
\hline
\end{tabular}
\vspace{0.35cm}
\begin{beamercolorbox}[sep=9pt,center,rounded=true]{block title}
\Large\bfseries Solution retenue : CORDIC hyperbolique
\end{beamercolorbox}
\end{frame}

% ============================================================
\section{L'algorithme CORDIC}

\begin{frame}{CORDIC : définition et principe général}
\begin{columns}[T]
\begin{column}{0.55\textwidth}
  \begin{block}{Définition}
    \small
    \textbf{CORDIC} (\textit{Coordinate Rotation Digital Computer}) calcule des fonctions trigonométriques et hyperboliques par une succession d'additions, de soustractions et de décalages binaires.
  \end{block}
  \begin{exampleblock}{Astuce matérielle}
    \[
      a\times 2^{-k}\equiv a \gg k
    \]
    \centering\small
    La multiplication par une puissance de deux devient un simple décalage à droite.
  \end{exampleblock}
\end{column}
\begin{column}{0.40\textwidth}
  \begin{block}{Deux modes}
    \begin{itemize}\small
      \item \textbf{Circulaire} : \(\sin\), \(\cos\).
      \item \textbf{Hyperbolique} : \(\sinh\), \(\cosh\), \(\tanh\).
    \end{itemize}
  \end{block}
  \begin{alertblock}{Relation utilisée}
    \[
      \sigma(x)=\frac{1+\tanh(x/2)}{2}
    \]
    \centering\small
    Un seul bloc hyperbolique permet donc de produire \(\tanh\) et \(\sigma\).
  \end{alertblock}
\end{column}
\end{columns}
\end{frame}

\begin{frame}{CORDIC circulaire \((m=+1)\) : \(\sin\) et \(\cos\)}
\begin{columns}[T]
\begin{column}{0.28\textwidth}
  \begin{block}{Équations itératives}
    \small
    \[
    \begin{cases}
      x_{k+1}=x_k-d_k2^{-k}y_k\\
      y_{k+1}=y_k+d_k2^{-k}x_k\\
      z_{k+1}=z_k-d_k\arctan(2^{-k})
    \end{cases}
    \]
    \[
      d_k=\mathrm{sign}(z_k)
    \]
  \end{block}
  \begin{exampleblock}{Résultat}
    \small
    Après \(n\) itérations : \(x_n\approx\cos(z_0)\), \(y_n\approx\sin(z_0)\).
  \end{exampleblock}
\end{column}
\begin{column}{0.70\textwidth}
 
  \includegraphics[width=\linewidth,height=5.5cm]{circul.png}
\end{column}
\end{columns}
\end{frame}

\begin{frame}{CORDIC hyperbolique \((m=-1)\) : \(\tanh\) et \(\sigma\)}
\begin{columns}[T]
\begin{column}{0.28\textwidth}
  \begin{block}{Équations itératives}
    \small
    \[
    \begin{cases}
      x_{k+1}=x_k+d_k2^{-k}y_k\\
      y_{k+1}=y_k+d_k2^{-k}x_k\\
      z_{k+1}=z_k-d_k\operatorname{atanh}(2^{-k})
    \end{cases}
    \]
  \end{block}
  \begin{exampleblock}{Sortie}
    \small
    \[
      x_n\approx \cosh(z_0),\qquad y_n\approx \sinh(z_0)
    \]
    \[
      \boxed{\tanh(z_0)=\frac{y_n}{x_n}}
    \]
  \end{exampleblock}
\end{column}
\begin{column}{0.70\textwidth}
 
  \includegraphics[width=\linewidth,height=5.5cm]{cord_hyper.png}
\end{column}
\end{columns}
\end{frame}

\begin{frame}{Pourquoi répéter certaines itérations en mode hyperbolique ?}
\begin{columns}[T]
\begin{column}{0.28\textwidth}
  \begin{block}{Problème}
    \small
    En mode hyperbolique, certaines corrections élémentaires ne suffisent pas à garantir la convergence sur toute la plage utile.
  \end{block}
  \begin{alertblock}{Solution}
    \small
    Les itérations critiques sont répétées :
    \[
      k=4,\;13,\;40,\;121,\ldots
    \]
    avec \(k_{i+1}=3k_i+1\).
  \end{alertblock}
  
\end{column}
\begin{column}{0.70\textwidth}

   \includegraphics[width=\linewidth,height=5.5cm]{iter.png}
   \begin{exampleblock}{Conséquence}
    \small
    Pour \(k=16\), les répétitions \(k=4\) et \(k=13\) conduisent à 18 passes.
  \end{exampleblock}
\end{column}
\end{columns}
\end{frame}

\begin{frame}{Choix du domaine de travail : influence sur le gradient de \(\tanh\)}
\begin{columns}[T]
\begin{column}{0.28\textwidth}
  \begin{block}{Contrainte}
    \small
    Le CORDIC hyperbolique converge naturellement sur un domaine limité :
    \[
      |z_0|<\sum_{k=1}^{n}\operatorname{atanh}(2^{-k})\approx 1.118.
    \]
  \end{block}
  \begin{exampleblock}{Choix pratique}
    \small
    Pour traiter \([-2,+2]\), on normalise l'entrée :
    \[
      z'=z/2 \in [-1,+1].
    \]
    Ce choix conserve une courbure utile sans saturation excessive.
  \end{exampleblock}
\end{column}
\begin{column}{0.70\textwidth}
  \figcol{interv.png}
  
\end{column}
\end{columns}
\end{frame}

\begin{frame}{Précision CORDIC en fonction du nombre d'itérations}
\begin{columns}[T]
\begin{column}{0.28\textwidth}
  \begin{block}{Règle pratique}
    \small
    Une itération CORDIC apporte approximativement un bit de précision supplémentaire.
  \end{block}
  \begin{exampleblock}{Choix de \(k=16\)}
    \small
    \begin{itemize}
      \item Cohérent avec une représentation fixe 16 bits.
      \item Erreur du même ordre que la résolution \(2^{-15}\).
      \item Compromis précision / latence / ressources.
    \end{itemize}
  \end{exampleblock}
\end{column}
\begin{column}{0.70\textwidth}

   \includegraphics[width=\linewidth,height=5.5cm]{preci.png}
\end{column}
\end{columns}
\end{frame}

\begin{frame}{Architecture CORDIC paramétrée : étages et flux de données}
\centering
\begin{tikzpicture}[x=1cm,y=1cm,
  bloc/.style={draw,rounded corners=3pt,minimum width=1.25cm,minimum height=0.72cm,font=\scriptsize\bfseries,align=center},
  arr/.style={-{Stealth[length=6pt]},very thick},
  brk/.style={decorate, decoration={brace,amplitude=5pt,raise=2pt}, thick}
]
  \node[bloc,fill=gray!20] (inp) at (0.4,2.7) {Entrée\\\(x\)};
  \node[bloc,fill=rougebox!25,draw=rougebox] (nm) at (2.0,2.7) {NP-RHC\\\(k\le0\)};
  \node[font=\large] (d1) at (3.3,2.7) {\(\cdots\)};
  \node[bloc,fill=bleubox!20,draw=bleubox] (p) at (4.8,2.7) {P-RHC\\\(1\ldots n\)};
  \node[bloc,fill=orangeiter!35,draw=orangeiter] (rep) at (6.6,2.7) {Rép.\\\(4,13\)};
  \node[bloc,fill=gray!30] (sum) at (8.1,2.7) {\(\Sigma\)};
  \node[bloc,fill=vertbox!20,draw=vertbox] (div) at (9.8,2.7) {VLC\\division};
  \node[bloc,fill=mygreen!20,draw=mygreen,minimum width=1.8cm] (out) at (11.8,2.7) {Sortie\\\(\tanh/\sigma\)};
  \foreach \a/\b in {inp/nm,nm/d1,d1/p,p/rep,rep/sum,sum/div,div/out} \draw[arr] (\a)--(\b);
  \draw[brk,rougebox] (1.25,3.35)--(2.75,3.35) node[midway,above=6pt,font=\tiny\bfseries\color{rougebox}]{plage};
  \draw[brk,bleubox] (4.05,3.35)--(7.35,3.35) node[midway,above=6pt,font=\tiny\bfseries\color{bleubox}]{précision hyperbolique};
  \draw[brk,vertbox] (9.1,3.35)--(10.55,3.35) node[midway,above=6pt,font=\tiny\bfseries\color{vertbox}]{division};

  \node[draw=rougebox,fill=rougebox!8,rounded corners,align=center,text width=3.4cm] at (2.1,0.95) {\textbf{Paramètre \(m\)}\\étend l'espace de convergence};
  \node[draw=bleubox,fill=bleubox!8,rounded corners,align=center,text width=3.6cm] at (5.9,0.95) {\textbf{Paramètre \(n\)}\\contrôle l'erreur principale};
  \node[draw=vertbox,fill=vertbox!8,rounded corners,align=center,text width=3.6cm] at (9.8,0.95) {\textbf{Paramètre \(p\)}\\affine la division finale};
  \draw[-{Stealth[length=5pt]},thick,dashed,rougebox] (2.1,1.42)--(2.1,2.25);
  \draw[-{Stealth[length=5pt]},thick,dashed,bleubox] (5.9,1.42)--(5.9,2.25);
  \draw[-{Stealth[length=5pt]},thick,dashed,vertbox] (9.8,1.42)--(9.8,2.25);
\end{tikzpicture}
\end{frame}

% ============================================================
%
\begin{comment}
 \section{Compléments expérimentaux}

\begin{frame}{Influence du réarrangement des données : shuffle}
\begin{columns}[T]
\begin{column}{0.28\textwidth}
  \begin{block}{Sans mélange}
    \small
    Les exemples sont présentés dans le même ordre ; l'erreur peut devenir cyclique et oscillante.
  \end{block}
  \begin{block}{Avec mélange}
    \small
    Le changement d'ordre à chaque époque réduit la corrélation entre mises à jour successives et stabilise la descente.
  \end{block}

\end{column}
\begin{column}{0.70\textwidth}

   \includegraphics[width=\linewidth,height=5.5cm]{shuffle.png}
\end{column}
\end{columns}
  \begin{exampleblock}{Interprétation}
    \small
    L'apprentissage du XOR et de la parité du Full Adder devient moins dépendant de l'ordre des exemples.
  \end{exampleblock}
\end{frame}

============================================================
\end{comment}
\section{Conclusion générale}

\begin{frame}{Conclusion}


\large
\begin{block}{Synthèse finale}
\begin{itemize}
    \item   Le Full Adder met en évidence la nécessité de la non-linéarité pour modéliser des fonctions complexes comme le XOR. 
    \item L’utilisation de \(\tanh\) permet l’apprentissage via un réseau multicouche, tandis que le seuillage final restitue la logique binaire. 
    \item Enfin, l’algorithme CORDIC offre une implémentation matérielle efficace de cette activation sur FPGA.
\end{itemize}
  




\end{block}


\end{frame}

\begin{frame}[plain]
\vspace{2.2cm}
\centering
{\Huge Merci pour votre attention}\\[0.5cm]
{\large Questions ?}
\end{frame}

\end{document}
